% Experiment report. Fixed sections: see experiments/TEMPLATE/report.tex and experiments/README.md.
% Build: npm run experiment:report -- 25-rl-spend-vi-r15
\documentclass{../polyexp}
\begin{document}
% generated by scripts/experiment-report.ts from experiment.json and models/<run>/ - do not edit
\providecommand{\ExpHeader}{}\renewcommand{\ExpHeader}{%
{\sffamily\footnotesize\color{expgrey}EXPERIMENT 25 \textperiodcentered\ METHOD \hfill 2026-11-27}\par
{\LARGE\bfseries\sffamily Spending actions and value iteration on the research-corrected base bot r15}\hfill\verdictbox{expred}{NOT KEPT}\par\smallskip
{\itshape With the base bot's research weight corrected (greedy-full-probe-perf-r15) and the new opponents, does value iteration over the spending actions still find a policy that beats the bot it steers?}\par\medskip
\begin{tabular}{@{}p{22mm}p{\dimexpr\linewidth-26mm}@{}}
\fact{Change}{Experiment 23 unchanged except the base bot (greedy-full-probe-perf-r15) and the opponents (rl-gate defaults after \#204: r15, probe-perf, anchor probe-perf).}
\fact{Knobs}{\knob{state} $\in$ \{cities,income,units,techs, turn,cities,income,units,techs\}, best turn,cities,income,units,techs; \knob{minCount} $\in$ \{1, 5, 20, 50\}, best 50}
\fact{Base}{\model{greedy-full-probe-perf-r15} $\rightarrow$ result \model{rl-gate vs greedy-full-probe-perf-r15: with turn FAIL (-240, 45.1\% better), without turn FAIL (-445); no fixed action beats r15 in the check}}
\fact{Method}{300k games with a random spend3 action per turn; value iteration for both state versions and min-count 1, 5, 20, 50; min-count chosen per version on 2,000 validation seeds; both chosen policies on the 10,000 evaluation seeds; metric: final raw score, \#171 rule}
\fact{Setups}{Perfection, 16$\times$16 continents, 4 players, Crazy, seat 0 Xin-Xi (human seat), learner base greedy-full-probe-perf-r15; opponents greedy-full-probe-perf-r15, greedy-full-probe-perf, greedy-full-probe-perf (anchor), seats shuffled per seed (rl-gate defaults from 27 Nov)}
\fact{Record}{issue \issue{166}; status: done}
\end{tabular}\par\medskip
\begin{center}\begin{tikzpicture}[node distance=5mm]
    \node[flowbox] (n0) {{\sffamily\bfseries\scriptsize\color{expgrey}BASE}\\[1pt]\model{full-probe-perf-r15}};
    \node[flowbox,right=5mm of n0] (n1) {{\sffamily\bfseries\scriptsize\color{expgrey}CHANGE}\\[1pt]\knob{state}\\{\scriptsize cities,income,units,techs, turn,cities,income,units,techs}\\\knob{minCount}\\{\scriptsize 1, 5, 20, 50}};
    \node[flowbox,right=5mm of n1] (n2) {{\sffamily\bfseries\scriptsize\color{expgrey}MEASURED}\\[1pt]Paired, 10,000 evaluation seeds\\{\scriptsize 1 setup, metric: final raw score, \#171 rule}};
    \node[flowbox,right=5mm of n2] (n3) {{\sffamily\bfseries\scriptsize\color{expgrey}RESULT}\\[1pt]Mean final score, 10,000 evaluation seeds: r15 vs spend3 policy (with turn)\\{\scriptsize 7199.0 $\rightarrow$ 6959.0}\\\textbf{$-$240.0}};
    \node[flowbox,right=5mm of n3] (n4) {{\sffamily\bfseries\scriptsize\color{expgrey}VERDICT}\\[1pt]\textbf{\color{expred}NOT KEPT}\\{\scriptsize new champion:}\\\model{rl-gate vs greedy-full-probe-perf-r15: with turn FAIL (-240, 45.1\% better), without turn FAIL (-445); no fixed action beats r15 in the check}};
    \draw[flowarrow] (n0) -- (n1); \draw[flowarrow] (n1) -- (n2); \draw[flowarrow] (n2) -- (n3); \draw[flowarrow] (n3) -- (n4);
\end{tikzpicture}\end{center}
\begin{center}\sffamily\small\begin{tabular}{@{}lrrr@{}}\toprule
Measure & Base & New & Change\\\midrule
Mean final score, 10,000 evaluation seeds: r15 vs spend3 policy (with turn) & 7199.0 & 6959.0 & $-$240.0\\
Same seeds: r15 vs spend3 policy without the turn & 7199.0 & 6754.0 & $-$445.0\\
\bottomrule\end{tabular}\end{center}
}
\providecommand{\ExpResults}{}\renewcommand{\ExpResults}{%
{\small\color{expgrey}No runs recorded.}
}

\ExpHeader

\section{Question and hypothesis}
\textbf{The idea (David, 27 Nov 2026, issue \#166).} Re-run experiment 23 with \texttt{greedy-full-probe-perf-r15}
as the base bot. In 23 the best policy (spending actions, state with the turn) gained +165 against probe-perf, but
probe-perf researched too little: raising its research weight ×1.5 alone (r15) gains +336 on the same gate. Most of
what value iteration found so far may have been this tuning gap. With r15 as the base, a gain would have to come
from deciding \emph{when} to spend on what.

\textbf{Prediction} (before any run of this experiment):
\begin{enumerate}
  \item \emph{Research every turn no longer helps.} In the action check, research ×2 on top of r15 (×3 against
        probe-perf, which gained less than ×1.5 there) is not better than r15: mean difference between $-150$ and
        $+50$. Economy every turn stays worse than the bot, army about even.
  \item \emph{The policy researches less} than in 23: research on fewer than 30\% of learner turns (23: 35.9\%).
  \item \emph{Little left to gain.} With the turn, the chosen policy's mean gain against r15 on the evaluation seeds
        is between $-50$ and $+100$ points, and it does \emph{not} pass the \#171 rule against r15.
  \item \emph{The turn still matters.} Without the turn the policy is again worse than the bot, and the sweeps do
        not settle.
\end{enumerate}
\textbf{Against the prediction:} if the policy passes against r15, value iteration finds timing decisions beyond
the research weight.

\section{Mechanism}
\textbf{Actions} (\texttt{spend3} in \texttt{rl\_value::strategic}): 0 economy (build), 1 army (train, attack),
2 research; each multiplies the base bot's move-type factors of its types by 2 for one learner turn (the same
boosts as these three actions in experiments 15--22). There is no ``bot unchanged'' action: a state without a policy
entry plays the base bot unchanged (David's choice). Army keeps train and attack (David's choice).

\textbf{State} (\texttt{--features}): (cities, income, units, techs), and (turn, cities, income, units, techs).
Dropped features are 0 in the policy table and in the lookup.

\textbf{Value iteration} as in experiment 21 ($Q(s,a) = $ mean of $r + V(s')$, $V(s) = \max_{a: n(s,a) \ge m}
Q(s,a)$, fallback to the mean over all transitions, $\gamma = 1$). With the turn, one backward sweep from the last
turn is exact. Without it, the same projected state occurs at different turns and can lead back to itself, so the
sweeps repeat (Jacobi, from $V = 0$) until the largest change is below $10^{-6}$, at most 30 times (the game's length
in learner turns); the tool reports the number of sweeps and the last change.

\section{Setup}
As experiment 23, with the base bot and the opponents of 27 Nov: learner base \texttt{greedy-full-probe-perf-r15};
opponents r15, probe-perf and the anchor probe-perf (\texttt{rl-gate} defaults after PR \#204), shuffled per seed.
Perfection, 16$\times$16 Continents, 4 players, Crazy, learner Xin-xi in the human seat. Seeds as in 23: training
$1\dots300{,}000$ (new games), action check $600{,}000{,}000\dots$ (2{,}000), validation $700{,}000{,}000\dots$
(2{,}000), evaluation $10^9\dots$ (10{,}000). Min-count chosen per state version on validation; both chosen policies
on the evaluation seeds with \texttt{rl-gate} against r15, and against the best fixed action of the check (if one
beats the bot) with \texttt{rl-value eval}.

Sanity check before the run: \texttt{rl-value} (default model now r15) and \texttt{rl-gate} give the same seat-0
score and win flag on 500 of 500 seeds (r15 against r15, seeds $1\dots500$).

Commands (from \texttt{engine/}): those of experiment 23 with \texttt{s3r15-} data files and
\texttt{../experiments/25-rl-spend-vi-r15} as output; \texttt{rl-gate --candidate
rl-policy:experiments/25-rl-spend-vi-r15/policy-<v>.tsv --base greedy-full-probe-perf-r15}.

\section{Results}
\ExpResults

All numbers are copied from \texttt{check.json}, \texttt{vi-summary-*.json}, \texttt{val-*.json} and the two
\texttt{rl-gate} JSONs in this folder (27 Nov 2026; per-seed rows dropped, the runs are deterministic). Data: 300{,}000
games with a random spend3 action per turn on r15, 8{,}449{,}051 rows, 1{,}485\,s on 10 threads.

\paragraph{Action check} (2{,}000 seeds; each action every turn, against r15):

\begin{center}\small
\begin{tabular}{lrrrrrl}
\toprule
action & mean & median & p10 & cities & techs & better / worse vs r15 \\
\midrule
bot (r15)  & 7148 & 8585 & 1510 & 2.62 & 19.31 & -- \\
economy    & 6687 & 7892 & 1365 & 2.48 & 18.34 & 629 / 1{,}051 \\
army       & 7157 & 8622 & 1500 & 2.62 & 19.36 & 418 / 427 \\
research   & 7112 & 8862 & 1424 & 2.32 & 19.96 & 997 / 922 \\
\bottomrule
\end{tabular}
\end{center}
No action beats the bot when played every turn (economy $-461$, army $+9$, research $-36$), so there is no fixed
action to compare the policy with on the evaluation seeds.

\paragraph{Value iteration and validation} (2{,}000 validation seeds, policy against r15, which scores 7{,}282 there;
start score 515):

\begin{center}\small
\begin{tabular}{lrrrrrrr}
\toprule
state & $m$ & covered & sweeps (last change) & optimism & policy mean & gain & better share \\
\midrule
no turn   & 1  & 100.0\% & 30 (755.0) & $+914$     & 6{,}739 & $-542$ & 40.42\% \\
no turn   & 5  & 98.9\%  & 30 (209.7) & $+743$     & 6{,}748 & $-534$ & 40.82\% \\
no turn   & 20 & 95.8\%  & 30 (197.6) & $+605$     & 6{,}760 & $-521$ & 41.40\% \\
no turn   & 50 & 91.3\%  & 30 (191.6) & $+506$     & 6{,}766 & $-515$ & \textbf{41.42\%} \\
with turn & 1  & 100.0\% & 1          & $+2{,}254$ & 6{,}971 & $-310$ & 46.19\% \\
with turn & 5  & 94.0\%  & 1          & $+1{,}348$ & 7{,}009 & $-273$ & 46.86\% \\
with turn & 20 & 81.5\%  & 1          & $+830$     & 7{,}033 & $-248$ & 46.89\% \\
with turn & 50 & 68.8\%  & 1          & $+524$     & 7{,}060 & $-221$ & \textbf{47.74\%} \\
\bottomrule
\end{tabular}
\end{center}
States: 53{,}521 without the turn, 258{,}897 with it. Min-count 50 is chosen for both versions.

\paragraph{Final test} (10{,}000 evaluation seeds $10^9\dots$, \texttt{rl-gate} against r15 with the default
opponents r15, probe-perf and the anchor probe-perf):

\begin{center}\small
\begin{tabular}{lrrrrl}
\toprule
policy & mean / p10 (r15) & better / worse & mean $d$ (99\% bounds) & 15 opponents (info) & rule \\
\midrule
with turn, $m=50$ & 6959 / 1435 (7199 / 1480) & 3947 / 4811 & $-240$ ($-295$, $-185$) & $-294$ & FAIL \\
no turn, $m=50$   & 6754 / 1425 (7199 / 1480) & 3707 / 5415 & $-445$ ($-502$, $-389$) & $-361$ & FAIL \\
\bottomrule
\end{tabular}
\end{center}
Both fail all three parts of the rule. Win rate 3.3\% (with turn) and 2.8\% (no turn) against r15's 3.8\%.

\section{Behaviour}
On the validation seeds the chosen policy with the turn plays economy 16.8\%, army 37.7\% and research 15.0\% of
learner turns, and the base bot unchanged (no policy entry) on 30.5\%; without the turn economy 44.7\%, army 35.4\%,
research 11.5\%. The fewer turns a policy takes an entry for (higher min-count), the less it loses: with the turn
$-310$ at $m=1$ (98.9\% of turns decided by the policy) up to $-221$ at $m=50$ (69.5\%).

\section{Interpretation}
\textbf{Prediction 1 (research every turn not better than r15, $-150$ to $+50$): right} ($-36$; economy and army as
predicted, $-461$ and $+9$).

\textbf{Prediction 2 (research on fewer than 30\% of turns): right} (24.4\% at $m=5$, 15.0\% at the chosen $m=50$).

\textbf{Prediction 3 (gain between $-50$ and $+100$, no PASS): no PASS right, range wrong.} The policy is clearly
worse than r15: $-240$ on the evaluation seeds.

\textbf{Prediction 4 (without the turn worse than the bot, sweeps do not settle): right} ($-445$; last change
192--755 after 30 sweeps).

\textbf{Why the policy is now worse than the bot.} With r15 as the base, the best choice in most turns is to leave the
bot as it is: none of the three boosts helps when played every turn. The spend3 action set has no ``bot unchanged''
action (David's choice in experiment 23, where the base still researched too little); the policy reaches the
unchanged bot only in states without an entry. So every decision the policy makes is a boost, and the noisy
maximum over three $Q$ estimates picks economy on a sixth of the turns, which costs 461 points when played every
turn. The more states fall back to the unchanged bot (higher min-count), the smaller the loss, which fits.

\textbf{What this says about experiments 21--23.} Their gains against probe-perf came from the research weight that
probe-perf lacked: r15 alone gains +336 there, 23's policy +165. Once the base bot is corrected, value iteration over
these one-turn boosts finds no timing decision that pays.

\todo{david}{Your reading: agree, disagree, what is missing?}

\section{Threats to validity}
\begin{itemize}
  \item \emph{The action set excludes the best single choice.} The result is about spend3 on r15; with ``bot
        unchanged'' as an action the policy could at least match the bot.
  \item \emph{Opponents changed with the base bot} (r15, probe-perf, anchor), so 25 and 23 are not paired.
\end{itemize}

\section{Verdict and next step}
\textbf{Not kept.} Both policies fail \texttt{rl-gate} against r15 ($-240$ with the turn, $-445$ without). David's RL v2
candidate stays experiment 23's policy, which was measured on the old base bot; against r15 it would be expected to
lose too.

Next step (one idea): put ``bot unchanged'' back into the action set (bot, economy, army, research) on r15, so the
policy has to justify every boost against doing nothing; together with a higher min-count this tests whether
\emph{any} turn-dependent boost pays on the corrected bot.

\section{Notes for the paper}
Claim: the gains of value iteration over strategic boosts (experiments 21--23) came from a mistuned base bot. After
correcting its research weight (r15: +336 against the old bot on its own), the same method gives a policy that loses
to the bot ($-240$), because none of the boosts pays when the base is well tuned and the action set had no ``do
nothing''. A clean example of measuring an RL gain against a weak baseline.

\todo{david}{Is this worth a place in the talk or report? Where?}

\end{document}
